\documentclass[11pt,a4paper]{article}
\usepackage[utf8]{inputenc}
\usepackage[T1]{fontenc}
\usepackage{amsmath,amssymb,mathrsfs}
\usepackage[margin=1in]{geometry}
\usepackage[colorlinks=true,linkcolor=blue,citecolor=blue,urlcolor=blue]{hyperref}

\title{The Ontology of Objective Time: Temporal Intensity, Diachronic Identity, and the Metaphysical Foundations of the \textsc{AEON}-$T_{\mu\nu}$ Framework}
\author{Lu\'is Pimenta Dam\'asio\thanks{Independent Researcher, Lisbon, Portugal. ORCID: \href{https://orcid.org}{0009-0001-7606-7557}. Master Compendium DOI: \href{https://doi.org}{10.5281/zenodo.24163322}. Observatory DOI: \href{https://doi.org}{10.5281/zenodo.24244501}.}}
\date{November 16, 2026}

\begin{document}
\maketitle

\begin{abstract}
Modern physics routinely treats spatial dimensions and external clock intervals as primary bedrock metrics, generating persistent conceptual paradoxes across quantum measurement, gravitational singularities, and cosmic entropy resets. Grounding inquiry in pre-mathematical ontology reverses this hierarchy: while space and duration are necessary postulates of existence, spatial lengths and dial readings are circumstantial observables bounded by local Operational Simultaneity Horizons ($\mathcal{L}_{\text{sim}}(r) > 0$). By demonstrating how conservation and interaction laws discipline an unbounded pre-mathematical infinity into a functional infinity of $100\%$-whole potentials, we define objective time as intrinsic temporal intensity ($\rho_\tau$). Formalizing diachronic identity $A_s^v$---distinguishing qualitative state $s$ from verification ordinal $v$---we prove logically that permanence is identical to change, every present is temporally composite and thick, and every coherent whole functions as an internal clock (\emph{Object-as-Subject}). Separating one-way ontic reason from the anachronic symmetry of formal mathematical equality (`='), we establish objective temporal intensity as the unifying metric across all physical domains.
\end{abstract}

\section{Introduction: Objective Time as the Unifying Metric}
Physics has long mistaken the meter reading for the pulse it measures, treating spatialized time axes ($dt \to 0$) and geometric distances as primary ontological building blocks. Pre-mathematical ontology restores the true order: spatial metrics and clock dials vary with horizon conditions, whereas intrinsic temporal intensity ($\rho_\tau$) and invariant Coherent Temporal Action ($h \cdot \mathcal{N}_{\text{eff}}$) remain invariant across every regime.

This constructive derivation stands on the shoulders of a rich philosophical lineage: St.~Augustine's non-metric lived present, Descartes' subjective starting point, Pascal's threshold of the infinite, Aristotle's lawful potentiality, and Leibniz's clock-bearing wholes. It is equally indebted to Kant's transcendental analysis of the subject--object relation, Hegel's sublation (\emph{Aufhebung}) of static identity, Nietzsche's dissolution of static substance into active becoming, Frege and Wittgenstein on the boundaries of formal identity (`='), and the thermodynamic irreversibility of Boltzmann and Prigogine.

\section{Pre-Mathematical Infinity, Functional Infinity, and Temporal Intensity}
Starting from pure subjective experience, time reveals an unbounded three-way infinity: no beginning, no end, and infinite divisibility between any two points. Space enters immediately by ontological necessity, because whatever lasts must exist somewhere. Even pre-mathematically, every existence obeys two foundational postulates: (1) \emph{if it exists, it has duration}; and (2) \emph{if it exists, it has a spatial coordinate}.

Yet unbounded infinity is completely homogeneous and cannot sustain metric distinctions. Setting the center of a circumference at infinity ($R = \infty$) flattens contiguous steps into a straight line and admits every random point on a two-dimensional plane as equidistant from the center. Conservation and interaction laws discipline this unboundedness into a \textbf{functional infinity}---an instant, immediate pool containing all lawfully permitted next states as $100\%$-whole potentials (a maximal entropic bundle). When a state instantiates out of this pool, it carries an intrinsic temporal signature prior to any clock reading, defining objective \textbf{temporal intensity ($\rho_\tau$)}.

\section{Diachronic Identity ($A_s^v$) and the Thick Present}
While \textbf{anachronic identity} ($A = A$) preserves the thread of continuity across states, real entities exist temporally and must be represented via \textbf{diachronic identity ($A_s^v$)}, where subscript $s$ denotes the qualitative state and superscript $v$ counts the sequential ordinal of verification of that state. From an initial state $A_1^1$, the very next moment splits into either \textbf{permanence} ($A_1^1 \to A_1^2$, where state $1$ endures and its verification ordinal advances) or \textbf{change} ($A_1^1 \to A_2^1$, where state $2$ emerges at verification ordinal $1$).

Because enduring requires active re-actualization across sequential verification ordinals, \textbf{permanence is ontologically identical to change}. Zero-width instants ($dt \to 0$) cannot exist: every actualized entity is temporally composite and occupies a \textbf{thick present} bounded physically by $\Delta E \cdot \Delta t_{\text{pre}} \ge \hbar/2$.

\section{Internal Clocks, Dual Planes, and Operational Horizons}
Because every coherent whole (``thing'') exists in a temporally composite moment, it obeys the \textbf{Second Postulate of Time}: \emph{if it exists, it is a clock}. Individuation operates across scales---from quantum excitations up to TON~618 across the \textsc{AEON}-$T_{\mu\nu}$ scale elevator ($10^{-35}\text{ m}$ to $10^{23}\text{ m}$)---with the Planck floor ($\Phi_{\text{floor}} = [1 + \pi(L/\ell_P)^2]^{-1}$) marking our projective horizon rather than an ontic void.

Whether viewed through radical empiricism or idealism, physical laws compel us to project an internal clock (\emph{Object-as-Subject}, $t_{\text{pre}}$) onto any phenomenal whole. Extensive clock duration ($t_{\text{obs}}$) and spatial rulers emerge only when an observer (\emph{Subject-as-Subject}) locks phase across an Operational Simultaneity Horizon ($\mathcal{L}_{\text{sim}}(r) = c\,n_{\text{eff}} / 2\Gamma_{\text{tot}} > 0$), bridged by invariant Coherent Temporal Action:
\begin{equation}
  h \cdot \mathcal{N}_{\text{eff}} \;=\; E_{\text{local}}\,t_{\text{pre}}\,\Phi_\tau(\Omega).
\end{equation}

\section{Present-Only Information and the Limits of Formal Equality (`=')}
Only the thick present is actual; matter stores no dead past, and every conceivable proof of past events is itself an active event of the now. What the ontic substrate provides is present structural information that, interpreted through conservation laws, enables lawful projection of presumed prior states (such as recognizing a fossil), while heat dissipation ($\Phi_\tau(\Omega)$) erases micro-correlations.

Retrodiction and prediction are therefore identical lawful operations applied to present structural states, leaving ultimate theological origins strictly open. Finally, formal mathematical equality (`=') operates in an anachronic, reversible plane ($t_{\text{obs}} \to -t_{\text{obs}}$), whereas real actualization obeys irreversible thermodynamic and \textbf{one-way ontic reason} ($A_1^1 \to A_1^2$ or $A_2^1$), establishing objective temporal intensity as the unifying metric across all physical fields.

\begin{thebibliography}{99}
\bibitem{aristotle} Aristotle, \emph{Physics} (Book IV, 217b--224a).
\bibitem{augustine} St.~Augustine, \emph{Confessions} (Book XI).
\bibitem{descartes} R.~Descartes (1641), \emph{Meditations on First Philosophy}.
\bibitem{pascal} B.~Pascal (1670), \emph{Pens\'ees}.
\bibitem{leibniz} G.~W.~Leibniz (1714), \emph{Monadology}.
\bibitem{kant} I.~Kant (1781/1787), \emph{Critique of Pure Reason} (Transcendental Aesthetic \& Dialectic).
\bibitem{hegel} G.~W.~F.~Hegel (1807), \emph{Phenomenology of Spirit}.
\bibitem{boltzmann} L.~Boltzmann (1877), ``\"Uber die Beziehung zwischen dem zweiten Hauptsatze der mechanischen W\"armetheorie und der Wahrscheinlichkeitsrechnung,'' \emph{Wiener Berichte} \textbf{76}, 373--435.
\bibitem{nietzsche} F.~Nietzsche (1887), \emph{On the Genealogy of Morality}.
\bibitem{frege} G.~Frege (1892), ``\"Uber Sinn und Bedeutung,'' \emph{Zeitschrift f\"ur Philosophie und philosophische Kritik} \textbf{100}, 25--50.
\bibitem{wittgenstein} L.~Wittgenstein (1921), \emph{Tractatus Logico-Philosophicus}.
\bibitem{prigogine} I.~Prigogine (1980), \emph{From Being to Becoming: Time and Complexity in the Physical Sciences} (W.~H.~Freeman).
\bibitem{damasio2026a} L.~Pimenta Dam\'asio (2026a), \emph{\textsc{AEON}-$T_{\mu\nu}$ \& Philosophy of Time Master Compendium}, Zenodo, DOI: \href{https://doi.org}{10.5281/zenodo.24163322}.
\bibitem{damasio2026b} L.~Pimenta Dam\'asio (2026b), \emph{\textsc{AEON}-$T_{\mu\nu}$: Covariant Kerr--Newman QFTCS \& Semiclassical Back-Reaction Observatory (\texttt{AEON-T14.html})}, Zenodo, DOI: \href{https://doi.org}{10.5281/zenodo.24244501}.
\bibitem{damasio2026c} L.~Pimenta Dam\'asio (2026c), \emph{Clockless Temporal Cadence, the Immediate Pool of Possibilities, and the Physical Mechanism of Conformal Cyclic Cosmology} (Paper 17), Zenodo, DOI: \href{https://doi.org}{10.5281/zenodo.24249370}.
\bibitem{damasio2026d} L.~Pimenta Dam\'asio (2026d), \emph{The Topography and Topology of Time: Coherent Temporal Action ($h \cdot \mathcal{N}_{\text{eff}}$), the Higgs Condensate, and Clockless Null Transit} (Paper 18), Zenodo, DOI: \href{https://doi.org}{10.5281/zenodo.24279981}.
\bibitem{damasio2026e} L.~Pimenta Dam\'asio (2026e), \emph{Relativity, Quantum Mechanics, Black Hole Singularities, and Quantum Gravity Under an Ontology of Intensive Time} (Paper 19), Zenodo, DOI: \href{https://doi.org}{10.5281/zenodo.24280496}.
\end{thebibliography}

\end{document}
